\begingroup
\renewcommand{\arraystretch}{1.45}
\Needspace{42\baselineskip}
\InterpolationArticle{37}

In Articles 30---36 we assumed that $\sin\mu a$ was not $=0$: it therefore remains to take up the same investigations for the case where $\sin\mu a=0$. Here we shall at once suppose that $a=0$ or $a=\frac{180^{\mathord{0}}}{\mu}$.
\SourceLeaf{329}
\par First let $X$ be a function of the form
\[
\begin{aligned}
&\gamma+\gamma'\cos x+\gamma''\cos2x+\mathrm{etc.}+\gamma^n\cos nx
\end{aligned}
\]
so that the number of unknown coefficients is $=1+n$. We already know that, among all the proposed values of the function $X$, those corresponding to a value of the argument greater than $180^{\mathord{0}}$ must be rejected: thus four cases must be distinguished here:

1) When $\mu$ is even and $a=0$, then $180^{\mathord{0}}$ will be the $\frac12\mu+1^{\mathrm{th}}$ value of $x$; hence, since the following ones must be rejected, $\frac12\mu+1$ values remain. Thus $n=\frac12\mu$ must hold.

2) Let $\mu$ be even and $a=\frac{180^{\mathord{0}}}{\mu}$. The $\frac12\mu^{\mathrm{th}}$ value of $x$ will be $180^{\mathord{0}}-\frac{180^{\mathord{0}}}{\mu}$; the following ones, which become greater than $180^{\mathord{0}}$, must be rejected. Thus $n=\frac12\mu-1$ must hold.

3) When $\mu$ is odd and $a=0$, for the value there results $\frac12\mu+\frac12^{\mathrm{th}}=180^{\mathord{0}}-\frac{180^{\mathord{0}}}{\mu}$, and

4) When $\mu$ is odd and $a=\frac{180^{\mathord{0}}}{\mu}$, for the value there results $\frac12\mu+\frac12^{\mathrm{th}}=180^{\mathord{0}}$: the following ones must be rejected in both cases, so that $n=\frac12\mu-\frac12$.

\Needspace{32\baselineskip}
Now, since the method used in the preceding articles could not be applied without some complications to the present case, where part of the given values would first have to be cut off from the complete period, we prefer the following method.

If, by the rules of Articles 20, 22, a function of the form
\[
\begin{aligned}
&\gamma+\gamma'\cos x+\gamma''\cos2x+\mathrm{etc.}+\gamma^m\cos mx\\[5pt]
&{}+\delta'\sin x+\delta''\sin2x+\mathrm{etc.}+\delta^m\sin mx
\end{aligned}
\]
is sought that satisfies all the $\mu$ given values, and in which $m=\frac12\mu-\frac12$ or $=\frac12\mu$ according as $\mu$ is odd or even, the coefficients $\delta'$, $\delta''$, $\delta'''$ etc. will of themselves become $=0$. For it is readily clear that, in an expression such as
\[
\begin{aligned}
&A\sin\lambda a+B\sin\lambda b+C\sin\lambda c+D\sin\lambda d+\mathrm{etc.}
\end{aligned}
\]
either the first part becomes $=0$, the last $=-B\sin\lambda b$, the next to last $=-C\sin\lambda c$, the one before that $=-D\sin\lambda d$ etc., namely when $a=0$; or the last becomes $=-A\sin\lambda a$, the next to last $=-B\sin\lambda b$, the one before that $=-C\sin\lambda c$ etc., when $a=\frac{180^{\mathord{0}}}{\mu}$, since for values of $x$ of which one is the complement of the other to $360^{\mathord{0}}$, the values of the function $X$ are equal. \Needspace{44\baselineskip}
Hence the function
\[
\begin{aligned}
&X'=\gamma+\gamma'\cos x+\gamma''\cos2x+\mathrm{etc,}+\gamma^m\cos mx
\end{aligned}
\]
\SourceLeaf{330}
in which the coefficients $\gamma$, $\gamma'$, $\gamma''$ etc. are determined by the formulas
\[
\begin{aligned}
&\gamma = \frac{1}{\mu}(A+B+C+D+\mathrm{etc.})\\[5pt]
&\gamma' = \frac{2}{\mu}(A\cos a+B\cos b+C\cos c+D\cos d+\mathrm{etc.})\\[5pt]
&\gamma'' = \frac{2}{\mu}(A\cos 2a+B\cos 2b+C\cos 2c+D\cos 2d+\mathrm{etc.})
\end{aligned}
\]
etc., but the last, when $\mu$ is even and thus $m=\frac12\mu$, by this one
\[
\begin{aligned}
&\gamma^m=\frac{1}{\mu}(A\cos ma+B\cos mb+C\cos mc+D\cos md+\mathrm{etc.})
\end{aligned}
\]
will necessarily be identical with the function $X$, provided that the latter is not of a higher degree than we defined above. For in cases $3$ and $4$, $X$ is of the $\frac12\mu-\frac12{}^{\mathrm{th}}$ order, i.e. the same as $X'$, and therefore, by Article 24, identical with $X'$. In the first case $X$ is of the same order as $X'$, and in the second of no higher order; hence in both cases all terms at least up to the $\frac12\mu-1^{\mathrm{th}}$ order will agree in the two functions (Article 24). That the terms of the $\frac12\mu^{\mathrm{th}}$ order must also agree in these functions follows, by the same Article 24, from the fact that in $X$ the conditional equation $K\sin ma=L\cos ma$ holds; namely, $L=0$, and in the first case $\sin ma=0$, while in the second, where $X$ rises only to the order $\frac12\mu-1$, $K=0$. In the second case, moreover, $X'$ seems to be of a higher order than $X$; but in this case the term of the $\frac12\mu^{\mathrm{th}}$ order in $X'$ also vanishes, since
\[
\begin{aligned}
&\gamma^{\frac12\mu}=\frac{1}{\mu}(A\cos 90^{{\mathord{0}}}+B\cos 270^{{\mathord{0}}}+C\cos 450^{{\mathord{0}}}+D\cos 630^{{\mathord{0}}}+\mathrm{etc.})=0
\end{aligned}
\]
so that in this case too $X'$ is in fact of the $m-1^{\mathrm{th}}$ order, or the $\frac12\mu-1^{\mathrm{th}}$.

\Needspace{46\baselineskip}
\InterpolationArticle{38}

If the function $X$ does not stop at the term $\cos nx$, but continues further, denoting the following terms by $\alpha^{n+1}\cos(n+1)x+\alpha^{n+2}\cos(n+2)x+\mathrm{etc.}$, there will be, by Articles 21, 23,
\[
\begin{aligned}
&\gamma = \alpha\pm\alpha^{\mu}+\alpha^{2\mu}\pm\mathrm{etc.}\\[5pt]
&\gamma' = \alpha'\pm\alpha^{\mu-1}\pm\alpha^{\mu+1}+\alpha^{2\mu-1}+\alpha^{2\mu+1}\pm\mathrm{etc.}\\[5pt]
&\gamma'' = \alpha''\pm\alpha^{\mu-2}\pm\alpha^{\mu+2}+\alpha^{2\mu-2}+\alpha^{2\mu+2}\pm\mathrm{etc.}\\[5pt]
&\gamma''' = \alpha'''\pm\alpha^{\mu-3}\pm\alpha^{\mu+3}+\alpha^{2\mu-3}+\alpha^{2\mu+3}\pm\mathrm{etc.}
\end{aligned}
\]
and so on as far as the last coefficient $\gamma^m$ when $\mu$ is odd, or the next to last $\gamma^{m-1}$ when $\mu$ is even; the lower sign applies here whenever $a=\frac{180^{{\mathord{0}}}}{\mu}$, and thus
\SourceLeaf{331}
in cases $2$ and $4$, the upper in cases $1$ and $2$; finally, for the last coefficient in the first case there will be
\[
\begin{aligned}
&\gamma^m=\alpha^m+\alpha^{3m}+\alpha^{5m}+\mathrm{etc.}
\end{aligned}
\]
but in the second case $\gamma^m=0$. These equations show how far it may be permissible to neglect the difference between the functions $X$ and $X'$. This latter function will be the simplest among all those that satisfy the $\mu$ proposed values; all these will be contained in the form $X'+Y\sin(\frac12\mu x-\frac12\mu a)$, which in cases $1$ and $3$ reduces to $X'+Y\sin\frac12\mu x$, but in cases $2$ and $4$ to $X'+Y\cos\frac12\mu x$. Evidently, however, if this expression is to be of the same form as $X$, i.e. free of sines, $Y$ must be
\[
\begin{aligned}
&\text{in the first case of the form} \quad  g\sin x+g'\sin 2x+g''\sin 3x+\mathrm{etc.}\\[5pt]
&\text{in the second case of the form} \quad  g+g'\cos x+g''\cos 2x+\mathrm{etc.}\\[5pt]
&\text{in the third case of the form} \quad  g\sin\frac12 x+g'\sin\frac32 x+g''\sin\frac52 x+\mathrm{etc.}\\[5pt]
&\text{in the fourth case of the form} \quad  g\cos\frac12 x+g'\cos\frac32 x+g''\cos\frac52 x+\mathrm{etc.}
\end{aligned}
\]
\Needspace{32\baselineskip}
And more generally, if $X''$ denotes any function similar to $X$ that satisfies the $\mu$ proposed values, all forms of this kind will be contained in the formula $X''+Y\sin\frac12\mu x$ or $X''+Y\cos\frac12\mu x$, where $Y$ denotes an undetermined function of the form that we have just presented. This can be carried out so that the function results reduced to the order $\frac12\mu$, $\frac12\mu-1$, $\frac12\mu-\frac12$, $\frac12\mu-\frac12$, and will evidently be identical with $X'$. The general rule for reducing such a function $X''$ to $X'$ is easily derived from Article 24. For any term $L\cos\lambda x$ in $X''$ one must substitute in $X'$, setting $\lambda=k\mu\pm\lambda'$ so that $\lambda'$ is not greater than $\frac12\mu$, the term $\pm L\cos\lambda'x$, where the lower sign must be taken whenever both $a=\frac{180^{{\mathord{0}}}}{\mu}$ and $k$ is even, and the upper in the remaining cases; finally, whenever in the second case, i.e. for $a=\frac{180^{{\mathord{0}}}}{\mu}$ and an even value of $\mu$, $\lambda'=\frac12\mu$ results, one must at once put, for $L\cos\lambda x$, $0$ in $X'$, or neglect that term entirely.

\endgroup
